% Generated draft: requires mathematical and source review.
\documentclass[11pt]{article}
\usepackage[margin=25mm]{geometry}
\usepackage{fontspec}
\setmainfont{Latin Modern Roman}
\usepackage{amsmath,amssymb,mathtools,unicode-math}
\setmathfont{Latin Modern Math}
\usepackage{enumitem,needspace,xurl,hyperref}
\setlength{\parindent}{0pt}
\setlength{\parskip}{5pt}
\setlength{\emergencystretch}{3em}
\newcommand{\N}{\mathbb{N}}
\newcommand{\Z}{\mathbb{Z}}
\newcommand{\Q}{\mathbb{Q}}
\newcommand{\R}{\mathbb{R}}
\newcommand{\C}{\mathbb{C}}
\newcommand{\F}{\mathbb{F}}
\newcommand{\A}{\mathbb{A}}
\newcommand{\PP}{\mathbb{P}}
\newcommand{\E}{\mathbb{E}}
\newcommand{\eps}{\varepsilon}
\newcommand{\dd}{\,\mathrm{d}}
\begin{document}
\section*{UnsolvedMath ID: 800913}
\textbf{Source problem code:} catalogue entry
\subsection{Definitions and mathematical statement}
Let \(\mathbb B^n=\{z\in\C^n:\|z\|_2<1\}\), where \(n\) is a positive integer, and let \(\mathrm{PSH}(\mathbb B^n)\) denote the plurisubharmonic functions on this ball. For \(K\subseteq\mathbb B^n\), put
\[
H_K(z)=\sup\{u(z):u\in\mathrm{PSH}(\mathbb B^n),\ u\leq 1\text{ on }\mathbb B^n,\ u|_K\leq 0\},
\qquad
h_K(z)=\limsup_{\substack{\zeta\to z\\\zeta\in\mathbb B^n}}H_K(\zeta).
\]
Thus \(h_K\) is the upper semicontinuous regularization of the displayed envelope. Following the supplied definition, call \(K\) regular when \(h_K\) is continuous. Here continuity is interpreted on \(\mathbb B^n\), the stated domain of the function; the supplied excerpt does not specify a boundary extension or boundary continuity requirement.

Write \(\operatorname{cap}\) for the capacity intended by the source. Its definition, normalization, ambient domain, and interpretation on arbitrary subsets are absent from the supplied record. Subject to resolving these qualifications, the question is whether
\[
\forall K\subseteq\mathbb B^n\text{ with }\operatorname{cap}(K)>0,\quad
\forall\varepsilon>0,\quad
\exists V\subseteq K:\quad
h_V\text{ is continuous on }\mathbb B^n
\quad\text{and}\quad
\operatorname{cap}(V)\geq\operatorname{cap}(K)-\varepsilon.
\]
The source says ``any subset'' and supplies no compactness or measurability hypothesis on \(K\) or \(V\). No such hypothesis is added here. It also supplies no finiteness assumption on capacity, so the intended meaning of the inequality when \(\operatorname{cap}(K)=+\infty\) requires clarification. The ball is the natural ambient set indicated by the envelope; the excerpt does not separately state the range of dimensions or the ambient-set hypothesis on \(K\).

The introductory assertion that a supremum is achieved by \(h_K\) refers to an optimization problem not included in the record. Consequently, that attainment assertion cannot be formulated precisely from the supplied evidence.
\subsection{Short English statement}
Question 45, attributed in the supplied record to Bloom--Leveberg, asks whether every set of positive capacity contains a regular subset whose capacity is within any prescribed \(\varepsilon>0\) of its own. Regularity means continuity of the relative extremal function. The capacity and some domain conventions remain unspecified in the supplied excerpt.
\subsection{Sources}
\begin{itemize}
\item Supplied local UnsolvedMath registry record: permanent ID \(800913\), exact source code \texttt{catalogue entry}, upstream ID \(800913\), upstream revision \texttt{d371ee917ac9647fbb06cd6fc529deec6f3dda82}. The registry labels the problem open; this is a supplied status, not an independently verified mathematical or literature assessment.
\item Dinew, Guedj, and Zeriahi, \emph{Open Problems in Pluripotential Theory} (2015), Question 45. Supplied source URL: \url{https://arxiv.org/abs/1511.00705}. The record reports extraction from arXiv LaTeX and access on 2026-10-16; these are supplied provenance details, not a claim of independent access or verification.
\item Attribution/context supplied by the registry: Bloom--Leveberg. The registry lists Dinew, Guedj, and Zeriahi as proposers and 2015 as the proposed year. No separate Bloom--Leveberg publication is identified in the supplied record.
\end{itemize}
\end{document}
